\documentclass[12pt,a4paper]{article}

\usepackage[T1]{fontenc}
\usepackage[utf8]{inputenc}
\usepackage{geometry}
\usepackage{hyperref}
\usepackage{longtable}
\usepackage{enumitem}
\usepackage{listings}
\usepackage{xcolor}
\usepackage{url}

\geometry{margin=2.5cm}
\hypersetup{
    colorlinks=true,
    linkcolor=blue,
    urlcolor=blue
}

\lstset{
    basicstyle=\ttfamily\small,
    breaklines=true,
    frame=single,
    columns=fullflexible
}

\title{Sudoku Game\\Opis użytych narzędzi i konfiguracji}
\author{}
\date{\today}

\begin{document}
\maketitle
\tableofcontents
\newpage

\section{Cel dokumentu}

Dokument opisuje narzędzia użyte przy przygotowaniu projektu Sudoku Game:
kompilację programu, generowanie dokumentacji programistycznej oraz generowanie
diagramu klas. Informacje można wykorzystać do odtworzenia środowiska i
ponownego wygenerowania artefaktów projektu.

\section{CMake}

CMake służy do konfiguracji i budowania programu C++.
Projekt korzysta z pliku \texttt{CMakeLists.txt} w katalogu głównym.

Najważniejsze ustawienia:
\begin{itemize}[noitemsep]
    \item minimalna wersja CMake: \texttt{3.16},
    \item język projektu: C++,
    \item standard języka: C++17,
    \item wynikowy program: \texttt{Sudoku\_Game}.
\end{itemize}

Typowy build lokalny:

\begin{lstlisting}
cmake -S . -B build
cmake --build build
\end{lstlisting}

Build pliku \texttt{.exe} dla Windows z poziomu Linuksa:

\begin{lstlisting}
cmake -S . -B build/windows-x64 \
  -DCMAKE_SYSTEM_NAME=Windows \
  -DCMAKE_CXX_COMPILER=x86_64-w64-mingw32-g++ \
  -DCMAKE_BUILD_TYPE=Release

cmake --build build/windows-x64 --parallel
\end{lstlisting}

Wariant z linkowaniem statycznym bibliotek GCC i standardowej biblioteki C++:

\begin{lstlisting}
cmake -S . -B build/windows-x64-static \
  -DCMAKE_SYSTEM_NAME=Windows \
  -DCMAKE_CXX_COMPILER=x86_64-w64-mingw32-g++ \
  -DCMAKE_BUILD_TYPE=Release \
  "-DCMAKE_EXE_LINKER_FLAGS=-static -static-libgcc -static-libstdc++"

cmake --build build/windows-x64-static --parallel
\end{lstlisting}

\section{Doxygen}

Doxygen został użyty do wygenerowania dokumentacji programistycznej na podstawie
komentarzy w plikach źródłowych C++.

Konfiguracja znajduje się w pliku \texttt{config/Doxyfile}. Najważniejsze parametry:

\begin{longtable}{p{0.35\textwidth}p{0.55\textwidth}}
\textbf{Parametr} & \textbf{Znaczenie} \\
\hline
\texttt{PROJECT\_NAME = "Sudoku"} & Nazwa projektu w dokumentacji. \\
\texttt{INPUT = core main.cpp} & Pliki i katalogi analizowane przez Doxygen. \\
\texttt{RECURSIVE = YES} & Rekurencyjne przeglądanie katalogów wejściowych. \\
\texttt{EXTRACT\_ALL = YES} & Do dokumentacji trafiają również elementy bez pełnych komentarzy. \\
\texttt{GENERATE\_HTML = YES} & Generowanie dokumentacji HTML. \\
\texttt{HTML\_OUTPUT = html} & Katalog wynikowy HTML wewnątrz \texttt{docs/generated/doxygen}. \\
\texttt{GENERATE\_LATEX = YES} & Generowanie dokumentacji LaTeX. \\
\texttt{LATEX\_OUTPUT = latex} & Katalog wynikowy LaTeX wewnątrz \texttt{docs/generated/doxygen}. \\
\texttt{HAVE\_DOT = YES} & Użycie Graphviz do diagramów zależności. \\
\texttt{CALL\_GRAPH = YES} & Generowanie grafów wywołań funkcji. \\
\texttt{CALLER\_GRAPH = YES} & Generowanie grafów funkcji wywołujących. \\
\texttt{DOT\_IMAGE\_FORMAT = png} & Format grafik Graphviz. \\
\end{longtable}

Polecenie generujące dokumentację:

\begin{lstlisting}
doxygen config/Doxyfile
\end{lstlisting}

Dokumentacja HTML jest dostępna w \texttt{docs/generated/doxygen/html/index.html}. Dokumentacja LaTeX
jest generowana w katalogu \texttt{docs/generated/doxygen/latex}. Plik PDF dokumentacji programistycznej
można zbudować poleceniem:

\begin{lstlisting}
make -C docs/generated/doxygen/latex
\end{lstlisting}

Wynikowy plik to \texttt{docs/generated/doxygen/latex/refman.pdf}.

\section{clang-uml}

\texttt{clang-uml} został użyty do wygenerowania diagramu klas z kodu C++.
Konfiguracja znajduje się w pliku \texttt{.clang-uml}.

Aktualna konfiguracja:

\begin{lstlisting}
compilation_database_dir: build
output_directory: docs/diagrams

diagrams:
  sudoku_class_diagram:
    type: class
    glob:
      - main.cpp
      - core/*.cpp
      - core/*.h
    include_system_headers: false
\end{lstlisting}

Znaczenie parametrów:
\begin{itemize}[noitemsep]
    \item \texttt{compilation\_database\_dir} wskazuje katalog z plikiem \texttt{compile\_commands.json};
    \item \texttt{output\_directory} wskazuje katalog wynikowy diagramów;
    \item \texttt{type: class} oznacza diagram klas;
    \item \texttt{glob} określa analizowane pliki projektu;
    \item \texttt{include\_system\_headers: false} pomija klasy z nagłówków systemowych.
\end{itemize}

Przed uruchomieniem \texttt{clang-uml} warto upewnić się, że istnieje aktualna
baza kompilacji:

\begin{lstlisting}
cmake -S . -B build -DCMAKE_EXPORT_COMPILE_COMMANDS=ON
\end{lstlisting}

Diagram generuje się poleceniem:

\begin{lstlisting}
clang-uml --query-driver /usr/bin/g++
\end{lstlisting}

Opcja \texttt{--query-driver /usr/bin/g++} pozwala \texttt{clang-uml} pobrać
ścieżki nagłówków systemowych z kompilatora użytego przez CMake. Bez tego na
niektórych instalacjach Clang może nie odnaleźć standardowych nagłówków C/C++.

Wyniki są zapisywane w katalogu \texttt{docs/diagrams}. Najważniejsze pliki:
\begin{itemize}[noitemsep]
    \item \texttt{docs/diagrams/sudoku\_class\_diagram.puml},
    \item \texttt{docs/diagrams/sudoku\_class\_diagram.svg},
    \item \texttt{docs/diagrams/sudoku\_class\_diagram.png}.
\end{itemize}

\section{LaTeX}

LaTeX został użyty do przygotowania dokumentacji przeznaczonej dla zwykłego
użytkownika oraz dokumentu opisującego użyte narzędzia. Źródła znajdują się w
katalogu \texttt{docs}.

Pliki źródłowe:
\begin{itemize}[noitemsep]
    \item \texttt{docs/user\_manual.tex} -- instrukcja użytkownika,
    \item \texttt{docs/tools\_used.tex} -- opis narzędzi i konfiguracji.
\end{itemize}

PDF-y można zbudować poleceniem:

\begin{lstlisting}
make -C docs
\end{lstlisting}

Wyniki:
\begin{itemize}[noitemsep]
    \item \texttt{docs/user\_manual.pdf},
    \item \texttt{docs/tools\_used.pdf}.
\end{itemize}

\section{Wersje narzędzi użyte podczas przygotowania dokumentacji}

W aktualnym środowisku użyto:
\begin{itemize}[noitemsep]
    \item Doxygen 1.16.1,
    \item clang-uml 0.6.2,
    \item CMake 4.3.3,
    \item x86\_64-w64-mingw32-g++ 16.1.0.
\end{itemize}

\end{document}
